\chapter{Hoarding}

\section*{Definition and Overview}
Hoarding is the persistent difficulty discarding or parting with possessions, regardless of their actual value, because of a perceived need to save them and distress at the thought of letting them go. When hoarding causes significant clutter that impairs living spaces, functioning, or safety, it is classified as hoarding disorder. Hoarding is more than collecting or being messy: it causes real distress, relationship strain, and health and safety hazards. It is a recognised mental health condition that responds to specialised psychological treatment and, where appropriate, support services.

\section*{Epidemiology}
Hoarding disorder affects around 2--6\% of the population, making it more common than previously recognised. It affects men and women in roughly equal numbers, though men may present more often to services because of fire and safety concerns. The condition typically begins in adolescence or early adulthood but becomes progressively more severe with age, most often coming to attention in people over 50. Many people with hoarding disorder also have depression, anxiety, or attention-deficit hyperactivity disorder.

\section*{Aetiology and Pathophysiology}
\begin{itemize}
    \item \textbf{Genetic factors:} hoarding runs in families, and twin studies show substantial heritability
    \item \textbf{Cognitive factors:} beliefs that items will be needed, have unique value, or hold memories
    \item \textbf{Emotional factors:} attachment to objects, difficulty making decisions, and avoidance of distress
    \item \textbf{Trauma and loss:} hoarding often worsens after losses, trauma, or significant life changes
    \item \textbf{Pathophysiology:} brain regions involved in decision-making, attachment, and impulse control function abnormally
    \item \textbf{Progression:} clutter accumulates gradually, and the person may have poor insight into the severity
    \item \textbf{Distinction from collecting:} collectors organise and display items with pleasure; hoarders are distressed and overwhelmed
\end{itemize}

\section*{Clinical Features}
\begin{itemize}
    \item Persistent difficulty discarding possessions regardless of value
    \item Strong urge to save items and distress at discarding them
    \item Clutter that fills rooms and prevents normal use of the home
    \item Living spaces, furniture, and appliances used for storage rather than their purpose
    \item Impaired functioning in work, social, and daily activities
    \item Relationship conflict with family members
    \item Health and safety hazards: fire risk, falls, infestation, and poor hygiene
    \item Distress, shame, and avoidance of visitors
\end{itemize}

\section*{Differential Diagnosis}
\begin{itemize}
    \item Normal collecting and disorganisation, which do not cause impairment
    \item Obsessive-compulsive disorder, where hoarding can occur but usually with obsessions and compulsions
    \item Depression, in which poor energy and motivation cause clutter
    \item Dementia and cognitive impairment causing disorganisation
    \item Psychosis, in which beliefs about possessions are delusional
    \item Attention-deficit hyperactivity disorder with chronic disorganisation
\end{itemize}

\section*{When to Seek Medical Care}
See a GP for assessment and referral when clutter affects safety, health, or relationships, or when the person or family is distressed. This is important before a crisis such as eviction, fire, or elder neglect occurs.

\textbf{Call 000 (or local emergency services) for:}
\begin{itemize}
    \item A fire in a cluttered home, or fire risk requiring urgent intervention
    \item A person at risk of harm from falls, collapse, or inability to access care
    \item Severe self-neglect or malnutrition
    \item Thoughts of self-harm or suicide
    \item An unsafe living situation involving a vulnerable person (also consider elder protection services)
\end{itemize}

\section*{Diagnosis and Investigations}
\begin{itemize}
    \item \textbf{Clinical interview:} assessment of difficulty discarding, distress, clutter, and impairment
    \item \textbf{Assessment scales:} standardised questionnaires measure hoarding severity
    \item \textbf{Home assessment:} where feasible, with consent, to assess clutter and safety
    \item \textbf{Mental health assessment:} exclude depression, OCD, ADHD, psychosis, and dementia
    \item \textbf{Functional assessment:} impact on work, health, and social life
    \item \textbf{Risk assessment:} fire, fall, and health risks are evaluated
\end{itemize}

\section*{Management}
\begin{itemize}
    \item \textbf{Cognitive behavioural therapy (CBT):} the main treatment, addressing beliefs, decision-making, and attachment to objects
    \item \textbf{Exposure and response prevention:} practising discarding while managing distress
    \item \textbf{Motivational interviewing:} builds readiness to change
    \item \textbf{Medication:} antidepressants, including SSRIs, may help coexisting depression and anxiety
    \item \textbf{Skills training:} organising, categorising, and decision-making skills
    \item \textbf{Multidisciplinary support:} occupational therapy, home help, and local council services
    \item \textbf{Family involvement:} education and support for families
    \item \textbf{Crisis management:} coordinated help for eviction, fire, and safety crises
\end{itemize}

\section*{Complications}
\begin{itemize}
    \item Fire risk and fire-related death, a leading cause of hoarding fatalities
    \item Falls and injuries in cluttered homes
    \item Infestations, mould, and poor sanitation affecting health
    \item Malnutrition and self-neglect
    \item Eviction, homelessness, and legal problems
    \item Social isolation and family breakdown
    \item Depression, anxiety, and reduced quality of life
\end{itemize}

\section*{Prognosis and Outcomes}
Hoarding disorder is treatable, though it tends to be chronic and requires persistence. CBT with a therapist experienced in hoarding produces meaningful improvement in many people, reducing clutter and distress. Relapse is common, and maintenance support helps sustain gains. Outcomes are best with early intervention, a strong therapeutic alliance, and a non-judgmental, collaborative approach. Families and services working together, rather than discarding possessions without consent, achieve the best and most lasting results.

\section*{Prevention and Self-Care}
\begin{itemize}
    \item Seek help early, before clutter becomes overwhelming
    \item Engage in CBT and practise discarding skills gradually
    \item Make small, regular decisions about possessions
    \item Use the "one in, one out" rule for new items
    \item Involve a trusted friend, family member, or support worker
    \item Treat coexisting depression and anxiety
    \item Attend support groups for people with hoarding problems
    \item Address safety hazards such as blocked exits and heating risks
\end{itemize}

\section*{Sources and Further Reading}
The following reputable sources provide further information for healthcare students and patients seeking reliable, current advice about hoarding.

\begin{itemize}
    \item SANE Australia. \textit{Hoarding disorder information and support.} https://www.sane.org/
    \item Healthdirect. \textit{Hoarding and hoarding disorder.} https://www.healthdirect.gov.au/hoarding
    \item Australian Psychological Society. \textit{Finding a psychologist for hoarding treatment.} https://psychology.org.au/
    \item NHS. \textit{Hoarding disorder: overview and treatment.} https://www.nhs.uk/mental-health/conditions/hoarding-disorder/
    \item Mayo Clinic. \textit{Hoarding disorder: symptoms and causes.} https://www.mayoclinic.org/
    \item International OCD Foundation. \textit{Hoarding center information.} https://hoarding.iocdf.org/
\end{itemize}
